\color{unidarkgreen}{\subsection{Логика закрытия отсечного клапана}}
\color{black}

\begin{enumerate}

\item
В некоторых случаях при срабатывании основных защит трансформатора и/или пуске пожаротушения трансформатора может потребоваться закрыть отсечной клапан. 

\item
Действие основных защит на отсечной клапан активируется отдельной накладкой \textbf{«Действие на ОК»}.

\item
При приеме сигналов срабатывания ДЗТ или ГЗ формируется команда \textit{«ЛЗ ОК: Закрытие»}. Длительность управляющего импульса задается уставкой \textbf{«Тимп ОК»}.

\item
Алгоритм работы ЛЗ ОК приведен на рисунке \ref{pic:look}. Параметры для настройки ЛЗ ОК приведены в таблице \ref{app_set:lzok}.

\medskip
\begin{figure}[H]
\centering
\includegraphics[width=1\textwidth,height=1\textheight,keepaspectratio]{\fbpath/06.01 ДЗТ/011.9902 ЛО ОК/img/look.pdf}
\captionof{figure}{Функциональная схема логики закрытия отсечного клапана}\label{pic:look}
\end{figure}


\end{enumerate}